\subsection{Rank-Nullity and Dimension Criteria}
\begin{theorem} \label{thm:range-spanned-by-images-of-basis} \hlt{Range from Basis Images}\\
Let $\mathcal{V}$ be finite-dimensional with basis $(\mathbf{v}_1,\ldots,\mathbf{v}_n)$, and let $T\in\LL(\mathcal{V},\mathcal{W})$. Then $\operatorname{range}(T)=\operatorname{span}(T\mathbf{v}_1,\ldots,T\mathbf{v}_n)$.\\
In particular, $\operatorname{range}(T)$ is finite-dimensional and $\dim(\operatorname{range}(T))\leq n$.
\begin{proof}
If $\mathbf{v}=c_1\mathbf{v}_1+\cdots+c_n\mathbf{v}_n$, then $T\mathbf{v}=c_1T\mathbf{v}_1+\cdots+c_nT\mathbf{v}_n$.\\
Thus every vector in $\operatorname{range}(T)$ lies in $\operatorname{span}(T\mathbf{v}_1,\ldots,T\mathbf{v}_n)$, and the reverse inclusion is immediate from the definition of range. The dimension bound follows because this spanning list has length $n$.\\
\hlt{Structural Commentary}: The range is generated by the images of any basis of the domain.\\
\hlt{Key Insight}: Linear maps carry a basis-spanning formula directly into a range-spanning formula.
\end{proof}
\end{theorem}

\begin{definition} \label{def:rank-and-nullity-linear-map} \hlt{Rank and Nullity}\\
For a finite-dimensional domain, the rank of $T$ is $\dim(\operatorname{range}(T))$ and the nullity of $T$ is $\dim(\ker(T))$.
\end{definition}

\begin{theorem} \label{thm:rank-nullity} \hlt{Rank-Nullity Theorem}\\
If $\mathcal{V}$ is finite-dim and $T \in \LL(\mathcal{V}, \mathcal{W})$, then $\operatorname{range}(T)$ is finite-dim and $\dim(\mathcal{V}) = \dim(\ker(T)) + \dim(\operatorname{range}(T))$.
\begin{proof}
By Proposition \ref{prop:kernel-and-range-are-subspaces} and Corollary \ref{cor:finite-dimensional-spaces-and-subspaces-have-bases}, $\ker(T)$ has a basis $(\mathbf{u}_1,\ldots,\mathbf{u}_m)$.\\
Extend it to a basis $(\mathbf{u}_1,\ldots,\mathbf{u}_m,\mathbf{w}_1,\ldots,\mathbf{w}_r)$ of $\mathcal{V}$ by Theorem \ref{thm:extending-independent-lists}.\\
If $\mathbf{v}=\sum\limits_{i=1}^m a_i\mathbf{u}_i+\sum\limits_{j=1}^r b_j\mathbf{w}_j$, then $T\mathbf{v}=\sum\limits_{j=1}^r b_jT\mathbf{w}_j$, because $T\mathbf{u}_i=\mathbf{0}$. Hence $(T\mathbf{w}_1,\ldots,T\mathbf{w}_r)$ spans $\operatorname{range}(T)$.\\
If $\sum\limits_{j=1}^r c_jT\mathbf{w}_j=\mathbf{0}$, then $\sum\limits_{j=1}^r c_j\mathbf{w}_j\in\ker(T)$, so for some scalars $d_i$, $\sum\limits_{j=1}^r c_j\mathbf{w}_j=\sum\limits_{i=1}^m d_i\mathbf{u}_i$.\\
Thus $\sum\limits_{j=1}^r c_j\mathbf{w}_j-\sum\limits_{i=1}^m d_i\mathbf{u}_i=\mathbf{0}$, and linear independence of extended basis forces every $c_j$ and $d_i$ to be zero.\\
Hence $(T\mathbf{w}_1,\ldots,T\mathbf{w}_r)$ is a basis of $\operatorname{range}(T)$. Thus $\dim(\ker(T))=m$, $\dim(\operatorname{range}(T))=r$, and $\dim(\mathcal{V})=m+r$.\\
\hlt{Structural Commentary}: An adapted basis splits $\mathcal{V}$ into kernel directions and directions detected by $T$.\\
\hlt{Key Insight}: Rank-nullity says every domain dimension is either killed by $T$ or appears in its range.
\end{proof}
\end{theorem}

\begin{corollary} \label{cor:map-to-smaller-space-not-injective} \hlt{Maps to Smaller Spaces Are Not Injective}\\
Let $\mathcal{V}$ and $\mathcal{W}$ be finite-dimensional vector spaces with $\dim(\mathcal{V})>\dim(\mathcal{W})$.\\
Then no linear map from $\mathcal{V}$ to $\mathcal{W}$ is injective.
\begin{proof}
Let $T\in\LL(\mathcal{V},\mathcal{W})$. By Theorem \ref{thm:rank-nullity}, $\dim(\ker(T))=\dim(\mathcal{V})-\dim(\operatorname{range}(T)) \geq \dim(\mathcal{V})-\dim(\mathcal{W})>0$.\\
Thus $\ker(T)$ contains a nonzero vector, so $T$ is not injective by Proposition \ref{prop:injective-iff-kernel-zero}.\\
\hlt{Structural Commentary}: Excess domain dimension forces collapse.\\
\hlt{Key Insight}: The dimension deficit must appear inside the kernel.
\end{proof}
\end{corollary}

\begin{corollary} \label{cor:map-to-larger-space-not-surjective} \hlt{Maps to Larger Spaces Are Not Surjective}\\
Let $\mathcal{V}$ and $\mathcal{W}$ be finite-dimensional vector spaces with $\dim(\mathcal{V})<\dim(\mathcal{W})$.\\
Then no linear map from $\mathcal{V}$ to $\mathcal{W}$ is surjective.
\begin{proof}
Let $T\in\LL(\mathcal{V},\mathcal{W})$. By Theorem \ref{thm:rank-nullity}, $\dim(\operatorname{range}(T))=\dim(\mathcal{V})-\dim(\ker(T))\leq\dim(\mathcal{V})<\dim(\mathcal{W})$.\\
Hence $\operatorname{range}(T)$ is a proper subspace of $\mathcal{W}$, so $T$ is not surjective.\\
\hlt{Structural Commentary}: A smaller domain cannot cover a larger codomain.\\
\hlt{Key Insight}: The range can never have dimension larger than the domain.
\end{proof}
\end{corollary}

\begin{corollary} \label{cor:equal-dimension-kernel-range} \hlt{Equal-Dimension Kernel-Range Criterion}\\
Let $\mathcal{V}$ and $\mathcal{W}$ be finite-dimensional vector spaces with $\dim(\mathcal{V})=\dim(\mathcal{W})$, and let $T\in\LL(\mathcal{V},\mathcal{W})$.\\
Then $\ker(T)=\{\mathbf{0}\}$ if and only if $\operatorname{range}(T)=\mathcal{W}$. Equivalently, $T$ is injective if and only if $T$ is surjective.
\begin{proof}
\begin{enumerate}[label=\roman*.]
\setlength{\itemsep}{0pt}
\item[$(\Rightarrow)$] If $\ker(T)=\{\mathbf{0}\}$, then rank-nullity gives $\dim(\operatorname{range}(T))=\dim(\mathcal{V})=\dim(\mathcal{W})$.\\
Since $\operatorname{range}(T)$ is a subspace of $\mathcal{W}$ with the same dimension, $\operatorname{range}(T)=\mathcal{W}$.
\item[$(\Leftarrow)$] If $\operatorname{range}(T)=\mathcal{W}$, then rank-nullity gives $\dim(\ker(T))=\dim(\mathcal{V})-\dim(\mathcal{W})=0$, hence $\ker(T)=\{\mathbf{0}\}$.
\end{enumerate}
\hlt{Structural Commentary}: Equal finite dimensions make injectivity and surjectivity equivalent.\\
\hlt{Key Insight}: In equal dimensions, having no kernel loss is exactly having full range.
\end{proof}
\end{corollary}

\begin{proposition} \label{prop:isomorphisms-preserve-bases} \hlt{Isomorphisms Preserve Bases}\\
Let $T:\mathcal{V}\to\mathcal{W}$ be an isomorphism and let $\beta=(\mathbf{v}_1,\ldots,\mathbf{v}_n)$ be a finite list in $\mathcal{V}$. Set $T\beta=(T\mathbf{v}_1,\ldots,T\mathbf{v}_n)$.\\
The following are equivalent:
\begin{enumerate}[label=(\roman*)]
\setlength{\itemsep}{0pt}
\item $\beta$ is a basis of $\mathcal{V}$.
\item $T\beta$ is a basis of $\mathcal{W}$.
\end{enumerate}
\begin{proof}
\begin{enumerate}[label=(\roman*).]
\setlength{\itemsep}{0pt}
\item[$(i)$]$\Rightarrow (ii)$ If $\mathbf{w}\in\mathcal{W}$, choose $\mathbf{v}\in\mathcal{V}$ with $\mathbf{w}=T\mathbf{v}$ and write $\mathbf{v}=\sum_j a_j\mathbf{v}_j$. Then $\mathbf{w}=\sum_j a_jT\mathbf{v}_j$, so $T\beta$ spans $\mathcal{W}$. If $\sum_j a_jT\mathbf{v}_j=\mathbf{0}$, applying $T^{-1}$ gives $\sum_j a_j\mathbf{v}_j=\mathbf{0}$; hence all $a_j=0$. Thus $T\beta$ is a basis.
\item[$(ii)$]$\Rightarrow (i)$
Assume $T\beta$ is a basis of $\mathcal{W}$. If $\mathbf{v}\in\mathcal{V}$, write $T\mathbf{v}=\sum_j a_jT\mathbf{v}_j$; applying $T^{-1}$ gives $\mathbf{v}=\sum_j a_j\mathbf{v}_j$, so $\beta$ spans $\mathcal{V}$. If $\sum_j a_j\mathbf{v}_j=\mathbf{0}$, applying $T$ gives $\sum_j a_jT\mathbf{v}_j=\mathbf{0}$; hence all $a_j=0$. Thus $\beta$ is linearly independent, and therefore a basis of $\mathcal{V}$.
\end{enumerate}
\hlt{Structural Commentary}: Isomorphisms preserve exactly the spanning and independence data that define a basis.\\
\hlt{Key Insight}: An isomorphism changes the ambient space, not the linear relations among vectors.
\end{proof}
\end{proposition}

\begin{theorem} \label{thm:finite-dimensional-operator-invertible-equivalences} \hlt{Finite-Dimensional Invertibility Criteria}\\
If $\mathcal{V}$ is finite-dimensional and $T \in \LL(\mathcal{V})$, then the following are equivalent:
\begin{enumerate}[label=\roman*.]
\setlength{\itemsep}{0pt}
\item $T$ is invertible;
\item $T$ is injective;
\item $T$ is surjective.
\end{enumerate}
\begin{proof}
\begin{enumerate}[label=(\roman)*.]
\setlength{\itemsep}{0pt}
\item[$(i)$] $\Rightarrow (ii)$ follows from Proposition \ref{prop:linear-map-invertible-iff-bijective}.
\item[$(ii)$] $\Rightarrow (iii)$: If $T$ is injective, then $\ker(T)=\{\mathbf{0}\}$.\\
Rank-nullity gives $\dim(\operatorname{range}(T))=\dim(\mathcal{V})$, hence $\operatorname{range}(T)=\mathcal{V}$, so $T$ is surjective.
\item[$(iii)$] $\Rightarrow (i)$: If $T$ is surjective, rank-nullity gives $\dim(\ker(T))=0$, so $T$ is injective.\\
Thus $T$ is bijective, hence invertible by Proposition \ref{prop:linear-map-invertible-iff-bijective}.
\end{enumerate}
\hlt{Structural Commentary}: For operators on finite-dimensional spaces, kernel and range dimensions trade off exactly. Injectivity and surjectivity therefore force each other.\\
\hlt{Key Insight}: In equal finite dimensions, losing no information is the same as reaching the whole space.
\end{proof}
\end{theorem}

\begin{perspective} \label{per:kernel-obstruction-to-invertibility} \hlt{Kernel Obstruction to Invertibility}\\
By Theorem \ref{thm:finite-dimensional-operator-invertible-equivalences} and Proposition \ref{prop:injective-iff-kernel-zero}, a finite-dimensional operator is invertible exactly when its kernel is zero. Any nonzero $\mathbf{v}\in\ker T$ therefore blocks invertibility.
\end{perspective}

\begin{theorem} \label{thm:finite-dimensional-isomorphic-iff-same-dimension} \hlt{Isomorphism and Dimension}\\
Two finite-dimensional vector spaces are isomorphic if and only if they have the same dimension.
\begin{proof}
\begin{enumerate}[label=\roman*.]
\setlength{\itemsep}{0pt}
\item[$(\Rightarrow)$] If $T:\mathcal{V}\to\mathcal{W}$ is an isomorphism and $(\mathbf{v}_1,\ldots,\mathbf{v}_n)$ is a basis of $\mathcal{V}$, then Proposition \ref{prop:isomorphisms-preserve-bases} shows that $(T\mathbf{v}_1,\ldots,T\mathbf{v}_n)$ is a basis of $\mathcal{W}$.
\item[$(\Leftarrow)$] If both spaces have dimension $n$, choose bases $(\mathbf{v}_1,\ldots,\mathbf{v}_n)$ and $(\mathbf{w}_1,\ldots,\mathbf{w}_n)$.\\
By Proposition \ref{prop:linear-map-determined-by-basis}, there is a linear map $T$ with $T\mathbf{v}_j=\mathbf{w}_j$ for every $j$. Then $T(\sum a_j\mathbf{v}_j)=\sum a_j\mathbf{w}_j$, so $\ker(T)=\{\mathbf{0}\}$ and $\operatorname{range}(T)=\mathcal{W}$. Hence $T$ is an isomorphism.
\end{enumerate}
\hlt{Structural Commentary}: Dimension classifies finite-dimensional spaces up to isomorphism.\\
\hlt{Key Insight}: Matching bases give an isomorphism; any isomorphism carries bases to bases.
\end{proof}
\end{theorem}
